In the Schwarzschild interior in synchronous time $\tau$ (with $\tau = 0$ the singularity), and equally in the Schwarzschild-de Sitter interior, whose metric functions differ by relative $O(\Lambda\tau^2)$ and for which the same limit holds with the same exponents and rates \cite[Thm.~5.9]{ButlerT6}, the orthonormal frame adapted to the Killing and sphere directions is parallel along $\partial_\tau$, and its connection is Kasner at leading order: the Hubble-normalised connection $\tau\,\omega$ tends to $\sum_i p_i\, B_i \otimes \hat\theta^i$ with $B_i$ the boosts, $\hat\theta^i$ the normalised coframe, and exponents $(p_1, p_2, p_3) = (-\tfrac13, \tfrac23, \tfrac23)$, with explicit rates: $O(\tau^{2/3})$ in the boost block and $O(\tau^{1/3})$ in the rotation block, the latter being the sphere's own connection, which is outside the Kasner comparison; the limit is a values-level statement (the algebra generated by the limiting connection values); the transport-level statement, convergence of the transport to the Kasner model's transport, holds along every tangential curve of fixed Hubble-normalised length, for every admissible metric of the geometry paper's class \cite[Def.~5.5]{ButlerT6}, whose exponents satisfy $\sum_i p_i^2 = 1$ \cite[Thm.~5.11]{ButlerT6}, and along the Schwarzschild $\varphi$-lines exactly, at rate $\tau^{1/3}$ \cite[Cor.~5.12]{ButlerT6}; what the geometry paper disclaims is the comparison over segments of fixed coordinate length, whose Hubble-normalised length diverges, and about which no claim is made here either. This coframe consists of sections of the sheaf paper's exceptional module, the weights of the sheaf paper's completion are forced to equal the exponents, and the total residue bivector $B_w = \sum_i p_i\, \hat e_0\wedge\hat e_i$ is a unit boost of the exceptional algebra. \emph{Provenance:} \cite[Lems.~5.1-5.4, Def.~5.5, Thm.~5.6, Thm.~5.9, Thm.~5.11, Cor.~5.12]{ButlerT6}; \cite[Thm.~8.16, Prop.~8.26, Thm.~8.27]{ButlerT4}.
